\section{Comunicação com os Envolvidos}

\paragraph{}
A comunicação entre as partes envolvidas no teste de penetração foi definida
com o objetivo de garantir alinhamento contínuo, transparência na execução das
atividades e rápida resposta a eventuais incidentes ou dúvidas relacionadas ao
escopo do projeto.

\subsection{Papéis e Responsabilidades}

\begin{itemize}
\item \textbf{Equipe de Pentest (Analistas):}
responsável pela execução dos testes, coleta de evidências, análise das
vulnerabilidades identificadas e elaboração dos relatórios técnicos.

\item \textbf{Gerente de Projeto ou Responsável pelo Laboratório:}
responsável pelo acompanhamento das atividades, aprovação de mudanças de escopo
e mediação entre a equipe técnica e os responsáveis pelo ambiente.

\item \textbf{Equipe Técnica do Laboratório:}
responsável por esclarecer dúvidas relacionadas à infraestrutura, aplicações e
serviços disponibilizados para avaliação.

\item \textbf{Equipe Jurídica:}
responsável pela validação do contrato formal de prestação dos serviços de
Pentest e do Acordo de Não Divulgação (NDA).

\item \textbf{Partes Interessadas (Stakeholders):}
devem ser informadas sobre vulnerabilidades críticas, alterações de escopo e
resultados relevantes identificados durante os testes.

\end{itemize}

\subsection{Canais de Comunicação}

Os canais oficiais definidos para comunicação durante o Pentest são:

\begin{itemize}
\item E-mail institucional;
\item Grupo oficial de comunicação da equipe (Teams, Slack ou equivalente);
\item Reuniões periódicas de acompanhamento;
\item Sistema de registro de chamados ou acompanhamento de atividades, quando
disponível.
\end{itemize}

\subsection{Janelas de Contato e SLA}

\begin{itemize}
\item Horário de atendimento: segunda a sexta-feira, das 08h às 18h;
\item Resposta para dúvidas operacionais: até 1 dia útil;
\item Resposta para incidentes críticos: até 4 horas;
\item Aprovação de mudanças de escopo: até 2 dias úteis;
\item Comunicação de vulnerabilidades críticas: imediata após validação da
evidência.
\end{itemize}

\subsection{Compartilhamento dos Resultados}

Com base na varredura de vulnerabilidades realizada com o OWASP ZAP, foi
elaborado um relatório resumido contendo as principais vulnerabilidades
identificadas, incluindo:

\begin{itemize}
\item SQL Injection;
\item Cross Site Scripting (Reflected);
\item Path Traversal;
\item External Redirect;
\item Ausência de proteção Anti-CSRF;
\item Ausência de cabeçalhos de segurança CSP e Anti-Clickjacking.
\end{itemize}

Os resultados deverão ser apresentados aos responsáveis pelo laboratório em
reunião de alinhamento, permitindo a discussão dos riscos identificados, dos
impactos potenciais e das próximas etapas do teste de penetração.

\subsection{Contrato e Acordo de NDA}

A equipe de Pentest colaborará na elaboração e revisão do contrato formal e do
Acordo de Não Divulgação (NDA), fornecendo informações técnicas necessárias
para delimitação do escopo, responsabilidades das partes e requisitos de
confidencialidade.

O NDA deverá proteger informações sensíveis obtidas durante os testes,
incluindo evidências técnicas, vulnerabilidades identificadas, credenciais,
informações de infraestrutura e quaisquer dados internos acessados durante a
execução das atividades.

\subsection{Gatilhos para Interrupção dos Testes (Stop-Test)}

A execução do Pentest deverá ser imediatamente interrompida e comunicada ao
responsável pelo laboratório caso ocorra alguma das seguintes situações:

\begin{itemize}
\item Indisponibilidade crítica dos sistemas avaliados;
\item Acesso acidental a dados pessoais reais ou informações sensíveis não
previstas no escopo;
\item Identificação de impacto operacional relevante ao ambiente;
\item Necessidade de alteração significativa do escopo aprovado;
\item Qualquer atividade que ultrapasse os limites previamente acordados.
\end{itemize}

\paragraph{}
Todas as comunicações relevantes deverão ser registradas e mantidas durante o
ciclo do projeto, garantindo rastreabilidade das decisões, aprovações e ações
executadas ao longo do teste de penetração.
